% ═══════════════════════════════════════════════════════════════
% LEHRERHINWEISE - DS-02: Mi nueva casa – Repaso
% Klasse 9, Unidad 2
% ═══════════════════════════════════════════════════════════════

\documentclass[11pt, a4paper]{article}

\usepackage[utf8]{inputenc}
\usepackage[T1]{fontenc}
\usepackage[ngerman]{babel}

\usepackage{helvet}
\renewcommand{\familydefault}{\sfdefault}

\usepackage[margin=2cm, top=2.5cm, bottom=2cm]{geometry}
\usepackage{enumitem}
\usepackage{tabularx}
\usepackage{booktabs}
\usepackage{array}
\usepackage{multicol}

\usepackage{fancyhdr}
\usepackage{lastpage}

\usepackage{xcolor}
\usepackage{tcolorbox}
\tcbuselibrary{skins}

% ═══════════════════════════════════════════════════════════════
% FARBEN
% ═══════════════════════════════════════════════════════════════

\definecolor{akzent}{HTML}{C41E3A}
\definecolor{hellgrau}{HTML}{F5F5F5}
\definecolor{mittelgrau}{HTML}{666666}
\definecolor{basis}{HTML}{4CAF50}
\definecolor{standard}{HTML}{2196F3}
\definecolor{erweiterung}{HTML}{9C27B0}

% ═══════════════════════════════════════════════════════════════
% BOXEN
% ═══════════════════════════════════════════════════════════════

\newtcolorbox{infobox}[1][]{
    colback=hellgrau, colframe=akzent, boxrule=1pt, arc=0pt,
    left=10pt, right=10pt, top=8pt, bottom=8pt,
    title={#1}, fonttitle=\bfseries, coltitle=white, colbacktitle=akzent
}

\newtcolorbox{diffbox}[2][]{
    colback=white, colframe=#2, boxrule=1pt, arc=0pt,
    left=8pt, right=8pt, top=6pt, bottom=6pt,
    title={#1}, fonttitle=\bfseries\small, coltitle=white, colbacktitle=#2
}

% ═══════════════════════════════════════════════════════════════
% KOPF- UND FUSSZEILE
% ═══════════════════════════════════════════════════════════════

\pagestyle{fancy}
\fancyhf{}
\renewcommand{\headrulewidth}{0.5pt}
\renewcommand{\footrulewidth}{0pt}
\lhead{\footnotesize\textsc{Lehrerhinweise} · DS-02}
\rhead{\footnotesize Mi nueva casa – Repaso · Klasse 9}
\cfoot{\footnotesize\thepage\,/\,\pageref{LastPage}}

\setlength{\parindent}{0pt}
\setlength{\parskip}{6pt}

% ═══════════════════════════════════════════════════════════════
% DOKUMENT
% ═══════════════════════════════════════════════════════════════

\begin{document}

% ═══════════════════════════════════════════════════════════════
% TITEL
% ═══════════════════════════════════════════════════════════════

\begin{center}
    {\LARGE\bfseries Lehrerhinweise: Mi nueva casa – Repaso}\\[6pt]
    {\large DS-02 · Klasse 9 · Unidad 2}
\end{center}

\vspace{8pt}

% ═══════════════════════════════════════════════════════════════
% ÜBERBLICK
% ═══════════════════════════════════════════════════════════════

\begin{infobox}[Stundenübersicht]
\begin{tabular}{@{}ll@{}}
\textbf{Thema:} & Mi nueva casa -- Repaso + Pretérito perfecto \\
\textbf{Dauer:} & 90 Minuten (Doppelstunde) \\
\textbf{Materialien:} & AB-02, LB-02, LP-02, SLIDES-02, Beamer \\
\textbf{Voraussetzung:} & DS-01 (Vokabeln, Bewegungsverben, Objektpronomen) \\
\textbf{Quelle:} & Qué pasa 2, Unidad 2 + Wiederholung Unidad 1 (pret. perf.) \\
\end{tabular}
\end{infobox}

\vspace{8pt}

% ═══════════════════════════════════════════════════════════════
% FEINZIELE
% ═══════════════════════════════════════════════════════════════

\section*{Feinziele}

\begin{itemize}[leftmargin=1.5em]
    \item \textbf{FZ1:} SuS \textbf{rufen} Hausvokabeln und Bewegungsverben korrekt ab (AFB I)
    \item \textbf{FZ2:} SuS \textbf{ordnen} Bewegungsverben den Präpositionen zu (AFB I)
    \item \textbf{FZ3:} SuS \textbf{bilden} Sätze im pretérito perfecto mit Bewegungsverben (AFB II)
    \item \textbf{FZ4:} SuS \textbf{kombinieren} pretérito perfecto mit Objektpronomen (AFB II)
    \item \textbf{FZ5:} SuS \textbf{verfassen} einen Umzugsbericht im pretérito perfecto (AFB III)
\end{itemize}

% ═══════════════════════════════════════════════════════════════
% STUNDENVERLAUF
% ═══════════════════════════════════════════════════════════════

\section*{Stundenverlauf}

\begin{tabularx}{\textwidth}{@{}c|l|X|c|l@{}}
\textbf{Min} & \textbf{Phase} & \textbf{Aktivität} & \textbf{SF} & \textbf{Material} \\
\midrule
5 & Einstieg & Vokabel-Quiz: Bilder zeigen, SuS raten & PL & SLIDES 3 \\
\midrule
8 & Wiederholung I & AB Aufgabe 1 bearbeiten (Vokabeln beschriften) & EA & AB-02 A1 \\
& & Selbstkontrolle mit AB-01 Kasten V1 & & \\
\midrule
5 & Wiederholung II & Bewegungsverben + Präpositionen (Memory/Zuordnung) & PL & SLIDES 5 \\
\midrule
5 & & AB Aufgabe 2 bearbeiten & EA & AB-02 A2 \\
& & Lösung besprechen & & \\
\midrule
5 & Input I & Pretérito perfecto reaktivieren: haber-Konjugation & PL & SLIDES 6-7 \\
\midrule
5 & Input II & Partizipien: regelmäßig + unregelmäßig (puesto, hecho) & PL & SLIDES 7 \\
\midrule
6 & Sicherung & AB Kasten G1 ergänzen & EA & AB-02 G1 \\
\midrule
2 & & AB Kasten R1 lesen (Redemittel) & EA & AB-02 R1 \\
\midrule
5 & Vertiefung & Kombination: pret. perf. + Objektpronomen einführen & PL & SLIDES 9 \\
& & Beispiele gemeinsam: Lo he puesto, La he llevado... & & \\
\midrule
10 & Üben I & AB Aufgabe 3 (Lückentext) bearbeiten & EA & AB-02 A3 \\
\midrule
3 & & Lösung besprechen & PL & \\
\midrule
3 & Input III & Modell für Satzbildung zeigen & PL & SLIDES 10 \\
\midrule
7 & Üben II & AB Aufgabe 4 (Sätze bilden) bearbeiten & EA & AB-02 A4 \\
\midrule
2 & & Lösung besprechen & PL & \\
\midrule
3 & Vorbereitung & Textproduktion einführen: Struktur, Konnektoren & PL & SLIDES 11 \\
\midrule
10 & Produktion & AB Aufgabe 5 (Umzugsbericht schreiben) & EA & AB-02 A5 \\
\midrule
5 & Sicherung & Freiwilliges Vorlesen + Peer-Feedback & PL & \\
\bottomrule
\end{tabularx}

\textbf{Gesamtzeit:} 89 Minuten

% ═══════════════════════════════════════════════════════════════
% DIFFERENZIERUNG
% ═══════════════════════════════════════════════════════════════

\section*{Differenzierung}

\begin{multicols}{3}

\begin{diffbox}[{[B] Basis}]{basis}
\footnotesize
\begin{itemize}[leftmargin=1em, itemsep=2pt]
    \item AB-01 Kasten V1 als Hilfe
    \item Satzanfänge für A5
    \item 6 Sätze ausreichend
    \item Mehr Zeit bei A3
\end{itemize}
\end{diffbox}

\columnbreak

\begin{diffbox}[{[S] Standard}]{standard}
\footnotesize
\begin{itemize}[leftmargin=1em, itemsep=2pt]
    \item Alle Aufgaben vollständig
    \item Eigene Formulierungen
    \item 6-8 Sätze in A5
    \item Normale Zeitvorgaben
\end{itemize}
\end{diffbox}

\columnbreak

\begin{diffbox}[{[E] Erweiterung}]{erweiterung}
\footnotesize
\begin{itemize}[leftmargin=1em, itemsep=2pt]
    \item Mehr als 8 Sätze in A5
    \item Objektpronomen integrieren
    \item Komplexere Satzstrukturen
    \item Peer-Tutoring für [B]
\end{itemize}
\end{diffbox}

\end{multicols}

% ═══════════════════════════════════════════════════════════════
% LÖSUNGEN
% ═══════════════════════════════════════════════════════════════

\section*{Lösungen}

\textbf{Kasten G1 (Pretérito perfecto):}
\begin{itemize}[leftmargin=1.5em, itemsep=0pt]
    \item haber: \textbf{he, has, ha, hemos, habéis, han}
    \item Partizipien: habl\textbf{ado}, com\textbf{ido}, viv\textbf{ido}
    \item Unregelmäßig: poner → \textbf{puesto}, hacer → \textbf{hecho}
\end{itemize}

\textbf{Aufgabe 1:}
\begin{enumerate}[leftmargin=1.5em, itemsep=0pt]
    \item la cocina, 2. la cama, 3. el sillón, 4. la ducha,
    \item el balcón, 6. el espejo, 7. el armario, 8. el pasillo
\end{enumerate}

\textbf{Aufgabe 2:}
\begin{itemize}[leftmargin=1.5em, itemsep=0pt]
    \item 1-b, 2-a, 3-c, 4-b, 5-a, 6-b
\end{itemize}

\textbf{Aufgabe 3:}
\begin{itemize}[leftmargin=1.5em, itemsep=0pt]
    \item nos hemos mudado
    \item han bajado
    \item ha subido
    \item he puesto
    \item la / he puesto
    \item ha sacado
    \item lo / ha llevado
\end{itemize}

\textbf{Aufgabe 4:}
\begin{enumerate}[label=\alph*), leftmargin=1.5em, itemsep=0pt]
    \item Clara ha puesto la lámpara en el salón.
    \item Nosotros hemos subido los muebles al piso.
    \item Papá ha sacado la nevera del camión.
    \item Yo he llevado el televisor al dormitorio.
    \item Mamá ha puesto el cuadro en la pared.
\end{enumerate}

% ═══════════════════════════════════════════════════════════════
% TYPISCHE FEHLER
% ═══════════════════════════════════════════════════════════════

\section*{Typische Fehler}

\begin{tabularx}{\textwidth}{@{}lXX@{}}
\textbf{Fehler} & \textbf{Beispiel} & \textbf{Korrektur} \\
\midrule
Pronomen nach haber & *He \textit{la} puesto... & \textbf{La} he puesto... \\
Falsches Partizip & *He \textit{ponido}... & He \textbf{puesto}... \\
haber nicht konjugiert & *Yo \textit{ha} puesto... & Yo \textbf{he} puesto... \\
Präposition vergessen & *He llevado el televisor \textit{el} salón. & ...televisor \textbf{al} salón. \\
Perfekt statt Präsens & *Siempre \textit{he puesto}... & Siempre \textbf{pongo}... \\
\bottomrule
\end{tabularx}

% ═══════════════════════════════════════════════════════════════
% HINWEISE
% ═══════════════════════════════════════════════════════════════

\section*{Didaktische Hinweise}

\begin{itemize}[leftmargin=1.5em]
    \item \textbf{Vokabel-Quiz:} Schnelles Warm-up, SuS bei Unsicherheit auf AB-01 verweisen
    \item \textbf{Pretérito perfecto:} Aus U1 bekannt -- kurze Reaktivierung reicht
    \item \textbf{Pronomen-Stellung:} Häufigster Fehler! Mehrfach betonen: Pronomen VOR haber
    \item \textbf{Textproduktion:} Satzanfänge auf Tafel/SLIDES für [B]-Schüler
    \item \textbf{Zeitmanagement:} Falls knapp, A5 als HA aufgeben
\end{itemize}

% ═══════════════════════════════════════════════════════════════
% AUSBLICK
% ═══════════════════════════════════════════════════════════════

\section*{Ausblick}

\textbf{DS-03 (Lunes 19):} Klausurvorbereitung -- alle Themen U1 + U2 (Stationenlernen)

\textbf{Klassenarbeit:} 19./20. Januar -- U1 (pretérito perfecto) + U2 (casa, verbos de movimiento, pronombres)

\textbf{Hausaufgabe:} AB Aufgabe 5 fertigstellen (falls nicht in Stunde geschafft)

\end{document}
